{
    \setlength{\parindent}{0em}
    \par {\zihao{4}\bfseries 一、题目：\Title}
    \\
    \par {\zihao{4}\bfseries 二、指导教师对文献综述、开题报告、外文翻译的具体要求：}
    \vspace{0.8em}

    {\zihao{5}
    \par 1. 文献综述和开题报告的一般要求：
    \par （1）查阅国内外相关文献资料 10 篇以上，其中外文文献不少于 5 篇，近 5 年文献不少于 3 篇。在此基础上完成文献综述，正文不少于 4000 字。
    \par （2）完成与可微渲染有关的外文翻译，正文（中文）不少于 3500 字。
    \par （3）完成开题报告，正文不少于 4500 字，内容应做到问题明确、方案可行、技术路线清晰、进度安排合理。

    \vspace{0.6em}
    \par 2. 对文献综述的具体要求：
    \par （1）了解鱼眼相机的成像模型和现有的相机标定管线。
    \par （2）梳理三维重建相关方法的发展脉络，重点总结传统多视图几何方法、NeRF 类方法、3DGS 类方法以及面向鱼眼或广角输入的适配研究，阐述其国内外研究现状。
    \par （3）比较已有方法在重建质量、渲染效率、训练代价等方面的差异，说明基于鱼眼相机和 3DGS 进行室内场景重建的研究意义。

    \vspace{0.6em}
    \par 3. 对开题报告的具体要求：
    \par （1）分析当前用于三维重建的主要表示方式的种类与特点，重点说明 3DGS 及其面向鱼眼输入时的适配思路。
    \par （2）给出基于鱼眼相机和雷达的场景重建设计方案，明确视频采集、相机标定、数据预处理、3DGS 重建等核心环节，并选择合理的处理算法。
    \par （3）提出论文的预期目标，包括实现可复现的鱼眼室内场景重建流程，并在重建质量或渲染效率方面取得具有分析价值的实验结果。
    \par （4）安排好论文进度，对文献调研、方案设计、系统实现、实验分析、论文撰写与修改等阶段给出详细规划。
    \par （5）如实说明本人已开展的调研、数据采集、初步重建等实验工作，并总结已取得的阶段性结果。}
}

\mbox{} \vfill

\signature{指导教师（签名）}
% comment the line above and uncomment the line below if you want to set a signature with a specific date.
% \signaturewithdate{指导教师（签名）}{1897}{5}{21}
